\section{Generic finite-volume uncle gaplessness}
\label{sec:uncle_finite_gaplessness}

\begin{lemma}[Vanishing ground-space projection]
    \label{lem:uncle_trial_ground_projection_decay}
    \lean{MPSTensor.uncleTwoWallTrialSource_span_projection_tendsto_zero}
    \notready
    \uses{lem:uncle_ground_gram_limits,thm:uncle_trial_ground_overlap_decay,lem:uncle_periodic_sector_projection_bound}
    Let $\varphi_N$ be the actual two-wall trial on $6N+1$ sites.
    Under the normalized block assumptions, let $P_N$ project onto the
    span of the two periodic block states. Then
    $\|P_N\varphi_N\|^2\longrightarrow0$.
\end{lemma}
\begin{proof}
    The two periodic squared norms tend to one and their cross inner
    product tends to zero. Eventually the Gram bounds of
    Lemma~\ref{lem:uncle_periodic_sector_projection_bound} hold.
    That estimate bounds the projected squared norm by four times
    the sum of the squared ground-state overlaps, which tends to zero.
\end{proof}

\begin{lemma}[The asymptotic variational inequality]
    \label{lem:uncle_asymptotic_variational_inequality}
    \lean{MPSTensor.eventually_exists_pos_eigenvalue_uncleHamiltonianES_lt_of_trial_bounds}
    \notready
    \uses{cor:uncle_variational_projection_bound}
    Let $H_N$ be the actual periodic uncle Hamiltonian on $6N+1$
    sites, and let $Q_N$ project onto its kernel. Suppose a trial family
    satisfies, eventually,
    \[
        \|v_N\|^2\ge c(N+1)^2,\qquad
        \operatorname{Re}\langle H_Nv_N,v_N\rangle\le CN,
        \qquad \|Q_Nv_N\|^2\longrightarrow0,
    \]
    for positive constants $c,C$. For every $\varepsilon>0$,
    $H_N$ eventually has an eigenvalue $0<\lambda<\varepsilon$.
\end{lemma}
\begin{proof}
    Eventually $\|Q_Nv_N\|^2<c/2$, so
    $\|v_N\|^2-\|Q_Nv_N\|^2\ge(c/2)N^2$.
    Taking $N>2C/(\varepsilon c)$ makes the energy strictly smaller
    than $\varepsilon$ times this difference. Apply
    Lemma~\ref{cor:uncle_variational_projection_bound}.
\end{proof}

\begin{theorem}[Generic finite-volume gaplessness]
    \label{thm:uncle_generic_finite_volume_gaplessness}
    \lean{MPSTensor.uncleHamiltonianES_ae_eventually_exists_pos_eigenvalue_lt}
    \notready
    \uses{thm:uncle_periodic_ground_space,thm:uncle_trial_generic_quadratic_norm,thm:uncle_trial_linear_energy,lem:uncle_trial_ground_projection_decay,lem:uncle_asymptotic_variational_inequality}
    Let $A,B$ be two blocks whose transfer maps are unital and whose
    dual transfer maps have faithful stationary density matrices of
    trace one. For almost every pair of walls $(R,L)$, if
    \[
        \operatorname{span}\left\{
        \begin{pmatrix}A_i&R_i\\L_i&B_i\end{pmatrix}:i
        \right\}
        =\operatorname{Mat}_{a+b}(\mathbb C),
    \]
    then, for every $\varepsilon>0$, the actual periodic uncle
    Hamiltonian on $6N+1$ sites eventually has an eigenvalue
    $0<\lambda<\varepsilon$.
\end{theorem}
\begin{proof}
    If no pair of walls satisfies the stated full-span condition,
    the implication holds for every pair. Otherwise fix one such
    pair. It supplies diagonal injectivity, primitive transfer decay
    and strict mixed damping. The quadratic norm coefficients are
    simultaneously positive outside a null set in the pair space.
    For each pair outside that set which satisfies full span, the
    actual kernel is exactly the span of the two periodic block
    states. The derived norm, energy and projection estimates now
    satisfy Lemma~\ref{lem:uncle_asymptotic_variational_inequality}.
\end{proof}

The full-span condition is the standing injective-perturbation assumption
of \cite[proof of the gaplessness theorem]{FernandezGonzalez2015Uncle}.
The quantifier concerns pairs satisfying that condition; no assertion
that almost every arbitrary wall pair satisfies it is needed here.
This is the finite-volume spectral conclusion. Construction of an
infinite-volume operator and determination of its full spectrum are
separate questions.
